% ==STYLE-BEGIN==
\documentclass[UTF8,11pt]{ctexart}
\usepackage{amsmath,amssymb,amsthm}
\usepackage[a4paper,margin=2.4cm,top=2.3cm,bottom=2.3cm]{geometry}
\usepackage{fancyhdr}
\usepackage{enumitem}
\usepackage{hyperref}
\hypersetup{hidelinks}

% ---- 版式 ----
\linespread{1.15}
\setlength{\parskip}{0.35em}
\setlength{\parindent}{1.6em}
\setlist{nosep,leftmargin=2em}

% ---- 页眉 ----
\pagestyle{fancy}
\fancyhf{}
\fancyhead[L]{\small\coursename}
\fancyhead[R]{\small\thepage}
\renewcommand{\headrulewidth}{0.4pt}

% ---- 定理环境：中文用正体，标题加粗，按小节共用一条流水号 ----
\newtheoremstyle{cnthm}{0.7em}{0.7em}{\normalfont}{}{\bfseries}{}{0.5em}{}
\theoremstyle{cnthm}
\newtheorem{thm}{定理}[section]
\newtheorem{lem}[thm]{引理}
\newtheorem{prop}[thm]{命题}
\newtheorem{cor}[thm]{推论}
\newtheorem{defn}[thm]{定义}
\newtheorem{rem}[thm]{注}
\newtheorem{exa}[thm]{例}
\newtheorem{conj}[thm]{猜想}
\renewcommand{\proofname}{\bfseries 证明}
\providecommand{\coursename}{}
% ==STYLE-END==
\renewcommand{\coursename}{高等代数\,·\,1015}
\title{高等代数 1015}
\date{}
\begin{document}
\maketitle

高等代数(线性代数)
研究对象:线性方程组,矩阵,行列式 ; 线性空间,线性变换,特征值,特征向量 ; 双线性函数,欧氏空间,Jordan标准型 ; 一元多项式环(初步)
目的: 研究线性空间上的线性变换

\section{复数}

$\mathbb{R}$实数集 $\mathbb{Q}$有理数集 $\mathbb{Z}$整数集 $\mathbb{Z_+}$正整数集
记$i$是一个符号
$C = \{ a+ib | a,b \in \mathbb{R} \}$
例如$3+i \cdot 2,3 + i(-1),0+i \cdot 2 \ \in C$
规定$C$上的运算
规定$a_1 +ib_1 = a_2 + ib_2 \Leftrightarrow a_1 = a_2 \ b_1 = b_2$
$0 + i \cdot 0$简记为0 
如果$b = 0 , a+ib = a$
\noindent\textbf{1.} 加法 $a_1+ib_1 , a_2+ib_2 \in C$
\[
	(a_1 + i b_1) + (a_2 + i b_2) \triangleq (a_1+a_2) + i(b_1+b_2) \in C
\]
	设$z_1 = a_1+ib_1 ,z_2 =  a_2+ib_2$,若
\[
	z_1+z_2 = 0 +i \cdot 0 \triangleq 0
\]
	则称$z_1$是$z_2$的相反数,记作$(-z_1)$
	此时有$-z_1 = (-a_1) +i(-b_1)$
	而减法则定义为$z_1-z_2 \triangleq z_1 + (-z_2)$
\[
	(a_1+ib_1)-(a_2+ib_2) = (a_1+ib_1)+((-a_2)+i(-b_2)) = (a_1-a_2)+i(b_1-b_2)
\]
	性质1 : C上的加法有如下性质
\noindent\textbf{1.} 结合律 $(z_1+z_2)+z_3 = z_1 + (z_2+z_3)$
\noindent\textbf{2.} 交换律 $z_1+z_2 = z_2+z_1$
\noindent\textbf{3.} 有零元 $\forall z \in C , \exists 0 + i \cdot 0 \in C \ s.t. \ z+0=0+z=z$
\noindent\textbf{4.} 有负元 $\forall z \in C , \exists z' \in C , s.t. z+z' = 0$
	于是$(C,+)$是一个Abel群(可交换群)
	代数系统 = 集合 + 运算 +法则
\noindent\textbf{2.} 乘法 规定 $i^2 = -1$
	有$z_1 = a_1 + ib_1 , z_2 = a_2 +ib_2$,有
\[
	z_1z_2 \triangleq (a_1a_2-b_1b_2)+i(a_1b_2+a_2b_1)
\]
	例如$z \cdot 0 = 0 , z \cdot 1 = z$
	设$z_1 = a_1 + ib_1 , z_2 = a_2 +ib_2$,如果$zz_1 = 1$,则称$z_2$是$z_1$的倒数
	如果$z_1z_2 = 1$,则$a_1a_2-b_1b_2 = 1$且$a_1b_2+a_2b_1 = 0$
	于是$a_2 = \frac{a_1}{a_1^2+b_1^2} , b_2 = \frac{-b_1}{a_1^2+b_1^2}$
	易知.当$z_1 \neq 0$(即$a_1 \neq 0$且$b_1 \neq 0$)时 , $z_1$必有唯一的倒数用$z_1^{-1}$表示
	除法 若$z_2 \neq 0 , \frac{z_1}{z_2} \triangleq z_1z_2^{-1}$
	性质2 : C上的乘法具有如下性质
\noindent\textbf{5.} 结合律 $(z_1 \cdot z_2) \cdot z_3 = z_1 \cdot (z_2 \cdot z_3)$
\noindent\textbf{6.} 交换律 $z_1 \cdot z_2 = z_2 \cdot z_1$
\noindent\textbf{7.} 有单位元$ \forall z \in C , 1 \cdot z = z \cdot 1 = z$
\noindent\textbf{8.} 有逆元$ \forall z \in C -\{0\} , \exists z' \in C , \ s.t. \ z \cdot z' = 1$
	于是$(C -\{ 0 \} , \cdot)$是一个Abel群
\noindent\textbf{9.} 分配律 $\forall z_1 , z_2 , z_3 \in C$ 
	$(z_1+z_2)z_3 = z_1z_3+z_2z_3$
	$z_1(z_2+z_3) = z_1z_2+z_1z_3$
	
	定理$(c , + , \cdot)$是一个域(field)
	C称为复数域 , 用$\mathbb{C}$表示
	
	例如,若P是素数,$F_p = \{ \bar{0} , \bar{1} , \cdots , \bar{p-1} \}$
\noindent\textbf{1.} $\bar{a} + \bar{b} = \bar{a+b}$
\noindent\textbf{2.} $\bar{a} \cdot \bar{b} = \bar{ab}$

\section{复数的几何意义}

在平面上 , 点$P(x,y)$与复数$x+iy$一一对应
即$\mathbb{C} \to \mathbb{R}^2$是一个双射
于是,$\mathbb{C}$也可以称为复平面
x轴称作实轴,y轴称作虚轴
x称作实部,y称作虚部,y称作虚部系数
$\vec{OP}$与$x$轴夹角称为辐角
若$z = a+ib$,规定
\[
\bar{z} = a-ib
\]
是z的共轭复数
性质3 : 
\noindent\textbf{1.} $\bar{z_1}+\bar{z_2} = \bar{z_1+z_2]}$
\noindent\textbf{2.} $\bar{z_1}\bar{z_2} = \bar{z_1z_2}$
\noindent\textbf{3.} $\frac{\bar{z_1}}{\bar{z_2}} = \bar{\frac{z_1}{z_2}}$
\noindent\textbf{4.} $\frac{z_1}{z_2} = \frac{z_1\bar{z_2}}{z_2 \bar{z_2}} = \frac{z_1\bar{z_2}}{|z_2|^2}$
规定$|z| = |\vec{OP}| = \sqrt{x^2+y^2}$
$z \cdot \bar{z} = |z|^2$

复数的三角表示

\end{document}
